% =============================================================================
% Копирование кода из PDF без искажений.
%
% Проблема: при копировании листинга из PDF теряются отступы, дефис
% подменяется типографским, а длинные строки рвутся. Решение (перенесено из
% основного учебника personal/lessons/tex/styles/preamble.tex): перед
% вёрсткой читаем исходный файл и прикрепляем к каждой строке её точный
% текст через /ActualText (пакет accsupp). Внешний вид не меняется.
% Нумерация строк в листингах не используется, чтобы номера не попадали в
% копируемый текст.
% =============================================================================
\usepackage{accsupp}
\makeatletter
\def\ACCSUPP@bdc{\special{pdf:code \ACCSUPP@span BDC}}
\def\ACCSUPP@emc{\special{pdf:code EMC}}
\makeatother

\ExplSyntaxOn
\seq_new:N  \g__codeacc_hex_seq
\int_new:N  \l__codeacc_index_int
\int_new:N  \l__codeacc_pending_int
\ior_new:N  \g__codeacc_ior
\bool_new:N \g__codeacc_active_bool
\tl_new:N   \l__codeacc_line_tl
\str_new:N  \l__codeacc_hex_str

% Store one raw source line as UTF-16BE hex (or an empty entry for a line
% that will not produce any glyph on the page).
\cs_new_protected:Npn \__codeacc_store_line:n #1
  {
    \tl_set:Nx \l__codeacc_line_tl { \tl_trim_spaces:n {#1} }
    \tl_if_empty:NTF \l__codeacc_line_tl
      {
        \int_incr:N \l__codeacc_pending_int
        \seq_gput_right:Nn \g__codeacc_hex_seq { }
      }
      {
        \str_set_convert:Nnnn \l__codeacc_hex_str {#1} { } { utf16be/hex }
        \seq_gput_right:Nx \g__codeacc_hex_seq { \str_use:N \l__codeacc_hex_str }
        \int_zero:N \l__codeacc_pending_int
      }
  }

% Read the whole source file once, before \inputminted typesets it.
\cs_new_protected:Npn \__codeacc_load:n #1
  {
    \seq_gclear:N \g__codeacc_hex_seq
    \bool_gset_false:N \g__codeacc_active_bool
    \file_if_exist:nT {#1}
      {
        \int_zero:N \l__codeacc_pending_int
        \ior_open:Nn \g__codeacc_ior {#1}
        \ior_str_map_inline:Nn \g__codeacc_ior { \__codeacc_store_line:n {##1} }
        \ior_close:N \g__codeacc_ior
        \bool_gset_true:N \g__codeacc_active_bool
      }
  }

% fvextra calls \FancyVerbFormatLine more than once per line (once inside a
% measuring \sbox), so the source line is looked up through fancyvrb's own
% FancyVerbLine counter instead of a counter of our own.
\cs_new_protected:Npn \__codeacc_wrap_line:n #1
  {
    \bool_if:NTF \g__codeacc_active_bool
      {
        \int_set:Nn \l__codeacc_index_int { \value{FancyVerbLine} }
        \str_set:Nx \l__codeacc_hex_str
          { \seq_item:Nn \g__codeacc_hex_seq { \l__codeacc_index_int } }
        \str_if_empty:NTF \l__codeacc_hex_str
          { #1 }
          {
            \exp_args:Nx \BeginAccSupp
              { method=hex,unicode,ActualText= \str_use:N \l__codeacc_hex_str }
            #1
            \EndAccSupp { }
          }
      }
      { #1 }
  }

% Оборачиваем \inputminted: так исправление работает и для \vfile, и для
% фрагментов vcode (tcolorbox выводит их через \inputminted из временного
% файла).
\cs_set_eq:NN \__codeacc_orig_inputminted:w \inputminted
\RenewDocumentCommand \inputminted { O{} m m }
  {
    \__codeacc_load:n {#3}
    \group_begin:
      \cs_set_eq:NN \FancyVerbFormatLine \__codeacc_wrap_line:n
      \__codeacc_orig_inputminted:w [#1] {#2} {#3}
    \group_end:
  }
\ExplSyntaxOff
